\documentclass[10pt]{article}
\usepackage[margin=.45in]{geometry}
\usepackage[T1]{fontenc}\usepackage{lmodern,microtype}
\usepackage{amsmath,amssymb,booktabs,enumitem}
\setlist{nosep,leftmargin=1.45em}
\pagestyle{plain}
\newcommand{\diag}{\operatorname{diag}}
\begin{document}
\begin{center}
{\Large\bf Exact Diffusion Design for Maximally Collective Stable Turing Patterns}\\[2pt]
{\large in Binary-Complex Mass-Action Networks}\\[2pt]
{\normalsize Topology-wide all-spectrum localization and exponent-optimal stationary heterogeneity trade-offs}\\[5pt]
{\bf Concise theorem summary --- external audit draft}
\end{center}

\paragraph{Family/realization space.}
For every $m\ge3$, the binary-complex mass-action topology $\widehat N_m$ has species $X_1,\ldots,X_m,Z$ and reactions
\[
0\to X_1,\quad X_1+X_i\to X_1+X_{i+1}\ (2\le i\le m-2),
\]
\[
X_1+X_{m-1}\to2X_m,\quad 2X_m\to X_2,
\quad 2Z\rightleftarrows X_1+X_m.
\]
The chain range is empty for $m=3$. With $n=m+1$, there are $n$ species and $n+1$ reactions. The semipositive conservation vector is $c=(0,4,\ldots,4,2,1)^T$, $\operatorname{rank}\Gamma_m=m$, and
\[
\ker\Gamma_m=\operatorname{span}\{(\mathbf1_m,0,0),(0_m,1,1)\}.
\]
The exact two-dimensional kernel and the $m+2$ columns give the stated rank by
rank--nullity; the maximal minor on rows $X_1,X_3,\ldots,X_m,Z$ and columns
$R_2,\ldots,R_{m-2},R_a,R_b,R_+$ has determinant $4(-1)^m$.
Every positive steady flux is $(a\mathbf1_m,b,b)$, and every positive-equilibrium Jacobian is exactly $J_m(a,b,H)=A_m(a,b)H$ with $a,b>0$ and $H\succ0$ diagonal. Conversely every such triple is physically realized.

\paragraph{Topology-wide localization.}
For every positive realization and every principal species set $I$ with $0<|I|<m$, $J_{I,I}$ is Hurwitz. Every strongly connected diagonal block is either a negative singleton, one of the two stable long cycles, or a non-singleton principal block of the boundary triad $\{X_1,X_m,Z\}$. The triad satisfies the cubic Routh--Hurwitz inequalities coefficientwise.  At $b=2a$, the deleted edge $X_1\to X_m$ lies on neither long cycle and can only refine the boundary decomposition.
For $m=3$, the two cycles $\{X_1,X_2\}$ and $\{X_2,X_3\}$ and all remaining
principal subgraphs are checked directly; the chain classification is used
for $m\ge4$.  In contrast,
\[
(-1)^m\det J_{\{X_1,\ldots,X_m\},\{X_1,\ldots,X_m\}}
=-2a^{m-1}b\prod_{i=1}^m h_i<0.
\]
Hence, throughout every positive realization,
\[
\min\left\{|I|:\varnothing\ne I\subseteq[n],\
\alpha(J_{I,I})>0\right\}=m=n-1.
\]
Every smaller subsystem is stable against both real and complex instability.

\paragraph{Principal-minor diffusion-ray theorem.}
Let $n\ge2$ and $a_I=(-1)^{|I|}\det J_{I,I}$.  Assume $\det J=0$, $a_I>0$ for
$|I|\le n-2$, and $\sum_{|I|=n-1}a_I>0$ (in particular these
coefficient hypotheses follow from a simple zero and otherwise stable
spectrum together with stable lower principal blocks).  Then for $D\succ0$
\[
\det(sD-J)=s(\beta_1+\beta_2s+\cdots+\beta_ns^{n-1}),\qquad \beta_k>0\ (k\ge2).
\]
A nonzero stationary threshold exists iff $\beta_1<0$; it is unique and simple. Moreover $J-sD$ has a positive real eigenvalue exactly for $0<s<s_*$. This does not rule out nonreal wave instability outside that band.  The
network application is a one-bad-minor corollary, not a hypothesis of the
general theorem.  Here scalar simplicity is first proved for the dispersion
root; ordinary algebraic simplicity of zero for $J-s_*D$ follows separately
from $\chi_{s_*}'(0)>0$.

\paragraph{Exact network diffusion law and sharp linear contrast.}
For $D=\diag(d_1,\ldots,d_m,d_Z)\succ0$, the complete order-$(n-1)$ omission table has one negative term (omit $Z$), positive weight $16$ when an interior $X_j$ is omitted, and zero when $X_1$ or $X_m$ is omitted. Thus
\[
\boxed{d_Z>8h_Z\sum_{j=2}^{m-1}\frac{d_j}{h_j}}
\]
is necessary and sufficient for a stationary threshold on the ray $sD$. For every $H$ in the explicitly defined homogeneously stable domain, the
sharp nonattained contrast infimum is
\[
T(H)=8h_Z\sum_{j=2}^{m-1}h_j^{-1};
\]
at $H=I$ it is $8(m-2)$. Every stationary crossing of this topology obeys
$\chi_D\chi_H>8(m-2)$.  Here $\chi_H$ measures equilibrium
concentration-scale separation; it is not interpreted as a universal measure
of biological expense or implausibility.

Here $c^TJ=0$ and homogeneous stability on $c^\perp$ imply rank $n-1$, an
algebraically simple conservation zero, and a positive order-$(n-1)$
characteristic coefficient.  Thus the general diffusion-ray hypotheses hold.
The criterion is independent of $a,b$, but the threshold is
$s_*(a,b,H,D)$ and generally depends on them.

\newpage
\paragraph{Exact unit-equilibrium stable design.}
At $a=b=1$, $H=I$, write $K_i=91m-181-i$ and take
\[
d_1=\frac{23}{63},\quad d_i=K_i^{-1}\ (2\le i\le m-1),
\quad d_m=\frac17,\quad d_Z=\frac{16}{45}.
\]
Write the full critical vectors as
$r=(r_1,\ldots,r_m,r_Z)^T$ and
$\ell=(\ell_1,\ldots,\ell_m,\ell_Z)^T$.
The exact contrast is $23(91m-183)/63$. A 35-term homogeneous half-plane
certificate and a 77-term spatial certificate exclude every nonzero
closed-right-half-plane first-mode root.  Algebraic simplicity at the equality
point follows from
\[
 \left.\frac{d}{d\lambda}\det(\lambda I-A_m+D_m)\right|_{\lambda=0}
 =-\frac{163}{45}\ell^Tr>0.
\]
Thus the first-mode kernel and cokernel are one-dimensional; exact critical
vectors also give transversality $\eta_m>0$.
The homogeneous base case $m=3$ is checked from
$q_3=(\lambda+7)(\lambda^2+5\lambda+2)$; the 35-term root exclusion applies
for $m\ge4$.

On the fixed integrated-mass class, the quadratic field satisfies
\[
A_mw_0=-\tfrac14B(r,r),\quad c^Tw_0=0,
\qquad (A_m-4D_m)w_2=-\tfrac14B(r,r).
\]
The cubic coefficient is
\[
c_m=\frac{\ell^T\{B(r,w_0)+\tfrac12B(r,w_2)\}}{\ell^Tr}.
\]
This is the one-dimensional center-manifold normal-form coefficient in the
adjoint-projection coordinate; setting the reduced vector field to zero gives
the stationary Lyapunov--Schmidt equation.
Spatial reflection $\xi\mapsto\pi-\xi$ sends $r\cos\xi$ to its negative
and acts as $A\mapsto-A$ on an equivariant center manifold.  Hence the reduced
field is odd and the two nonzero branches are exchanged by reflection.
Because their perturbations tend to zero in $H_N^2\hookrightarrow C^0$ from
$\mathbf1$, both branches are componentwise strictly positive for sufficiently
small $\mu>0$ at each fixed $m$.
For the stationary onset linearization
$\mathcal L_0=A_m+D_m\partial_{\xi\xi}$, fixed-mass invariance follows from
$c^TA_m=0$ and the Neumann boundary conditions.  The cosine blocks are
$A_m-k^2D_m$: the $k=0$ block is invertible on $c^\perp$, $k=1$ has the
certified simple zero, and $k\ge2$ is invertible.  The factorization
\[
 A_m-k^2D_m=-k^2D_m(I-k^{-2}D_m^{-1}A_m),\qquad
 \|(A_m-k^2D_m)^{-1}\|=O(k^{-2}),
\]
maps compatible $L^2$ data into $H_N^2$ and proves that the stationary
fixed-mass linearization is Fredholm of index zero, with
kernel $\operatorname{span}\{r\cos\xi\}$, cokernel
$\operatorname{span}\{\ell\cos\xi\}$, and nonzero transversality pairing
$(\pi/2)\ell^TD_mr$.  A four-factor recurrence and a shifted-polynomial
certificate prove $c_m<0$ for every $m\ge3$. The two supercritical positive
branches are locally exponentially asymptotically stable in the
fixed-integrated-mass $H^1$ phase space.

\paragraph{Stable diffusion--equilibrium trade-off.}
Let $\nu=m-2$,
\[
\kappa_\nu=\begin{cases}1/\sqrt3,&\nu=1,\\ \sqrt5/2,&\nu\ge2,
\end{cases}
\qquad L_0=\kappa_\nu/\sqrt\nu,
\qquad L_1=90\nu/(90\nu+1),
\]
and for $L\in[L_0,L_1]$ set
\[
h_1=h_m=h_Z=1,\qquad h_i=\frac{K_i}{LK_{i-1}},\qquad
\mathsf H_m(L)=\diag(h),\qquad
D_m^{\rm phys}(L)=\mathsf H_m(L)\Delta_m,
\]
where $\Delta_m$ is the effective unit profile above. The physical equilibrium
is $\mathsf H_m(L)^{-1}\mathbf1$.  Under
$\widehat x=\mathsf H_m(L)x$, the scaled PDE retains the parameter as
$\partial_t\widehat x=\mathsf H_m(L)\{f_m(\widehat x)
+(1-\mu)\Delta_m\partial_{\xi\xi}\widehat x\}$; at $\mu=0$ its
stationary bracket and critical right vector are unchanged.  The interval is
nonempty in every dimension.  The transformed left vector
$\widetilde\ell(L)=\mathsf H_m(L)^{-1}\ell$ satisfies
\[
 \widetilde\ell(L)^T\mathsf H_m(L)\Delta_mr
 =\ell^T\Delta_mr<0,
 \qquad \widetilde\ell(L)^Tr<0,
\]
so the crossing coefficient is positive.  The kernel of the row-scaled
first-mode operator remains
$\operatorname{span}\{r\}$; the nonzero pairing
$\widetilde\ell(L)^Tr$ excludes a generalized eigenvector, so the scaled
zero is algebraically simple.  A
22-term homogeneous certificate for $\nu\ge2$, a direct Routh--Hurwitz
calculation for $\nu=1$, an 84-term spatial certificate, and a
gauge-corrected cubic comparison prove a primary supercritical branch that is
locally exponentially asymptotically stable in the fixed-integrated-mass
$H^1$ phase space for every certified $L$.  The normalized branch is positive
for small $\mu>0$ by $H_N^2\hookrightarrow C^0$, and the positive diagonal
physical rescaling preserves componentwise positivity.  The physical contrasts are
\[
\chi_D(L)=\frac{23}{63}\,91\nu L,
\qquad
\chi_H(L)=\frac{91\nu-1}{91\nu L},
\qquad
\chi_D\chi_H=\frac{23(91\nu-1)}{63}.
\]
The product is independent of $L$ and equals the unit-equilibrium
stable-profile contrast,
\[
 \chi_D(L)\chi_H(L)=\chi_D^{\rm unit}(m)
 =\frac{23(91m-183)}{63}.
\]
Thus the family reallocates a fixed contrast product without reducing it.
Throughout the certified interval, $\chi_D(L)>\chi_H(L)$, so
$\max(\chi_D,\chi_H)=\chi_D(L)$ is uniquely minimized within this family at
$L=L_0$.
At $L=L_0$,
\[
 \chi_D=\frac{2093}{63}\kappa_\nu\sqrt\nu,
 \qquad
 \chi_H=\frac{\sqrt\nu}{\kappa_\nu}\frac{91\nu-1}{91\nu},
\]
so both, and hence their maximum, are $\Theta(\sqrt m)$, compared with
$\Theta(m)$ for the unit-equilibrium stable profile.  The lower endpoint is the
boundary of the present sufficient homogeneous modulus estimate for
$\nu\ge2$; it is not a proved intrinsic dynamical boundary. Since every
stationary crossing of $\widehat N_m$ satisfies
$\max(\chi_D,\chi_H)>\sqrt{8(m-2)}$, exponent $1/2$ is optimal among
stationary-crossing realizations of this topology. No constant-optimal or
complete Pareto-frontier claim is made.

\paragraph{PDE conclusion.}
On fixed-mass $L^2$, positive diagonal Neumann diffusion with domain $H_N^2$
is sectorial with compact resolvent, and the reaction Jacobian is a bounded
multiplication perturbation.  The high-mode spectral bound tends to
$-\infty$, so only finitely many low-mode gaps require analytic continuation
along the branch.  The half-order phase space is $H^1$, and the quadratic
mass-action Nemytskii map is smooth into $L^2$.  The conservation restriction
removes the homogeneous zero and the Neumann interval has no continuous
translation mode.  Strict branch-spectrum negativity therefore yields local
exponential convergence for sufficiently close nonnegative initial data in
the same mass class.

\paragraph{Most failure-prone points.}
External checking should focus on: (i) exhaustion of principal SCCs, including
$b=2a$; (ii) strict equality cases in the 22/35/77/84-term certificates and
the direct $\nu=1$ cubic; (iii) the principal-minor derivative monotonicity;
(iv) the equilibrium-scaled fixed-mass gauge; and (v) the cubic comparison
$N_m(L)>1/200$.

\paragraph{Scope.}
The family is synthetic, rank deficient with one semipositive conservation law, and all-mobile. Results are local for each fixed dimension. No arbitrary-data global boundedness, global attraction, constant-optimal stable frontier, biochemical realization, cross-diffusion, or immobile-species theorem is claimed.
Here Turing bifurcation means a stationary diffusion-driven crossing of a
nonconstant Neumann mode on a fixed interval; intrinsic finite-wavelength
selection under domain enlargement is not claimed.
\end{document}
